% 编译：xelatex 实验报告.tex && xelatex 实验报告.tex
\documentclass[UTF8,a4paper,12pt]{ctexart}

\usepackage[margin=2.5cm]{geometry}
\usepackage{graphicx}
\usepackage{amsmath}
\usepackage{booktabs}
\usepackage{hyperref}
\usepackage{caption}
\usepackage{subcaption}
\usepackage{float}
\usepackage{adjustbox}
\usepackage[hyphens]{url}
\usepackage{xurl}

\urlstyle{same}
\def\UrlFont{\ttfamily}
\newcommand{\varname}[1]{\url{#1}}

\newcommand{\fitfigure}[1]{%
    \adjustbox{max width=\linewidth,max height=0.82\textheight,keepaspectratio}{%
        \includegraphics{#1}}%
}

\hypersetup{
    colorlinks=true,
    linkcolor=black,
    citecolor=blue,
    urlcolor=blue
}

\title{\textbf{家庭电力消耗多变量时间序列预测} \\[0.5em]
\large 2026年专硕机器学习课程项目实验报告}
\author{%
学号：{\underline{\makebox[5cm][c]{20255227075}}} \quad
姓名：{\underline{\makebox[3cm][c]{刘通}}} \\[0.8em]
代码仓库：\url{请填写 GitHub 链接}%
}
\date{}

\begin{document}

\maketitle
\thispagestyle{empty}
\vspace{1em}
\noindent\textbf{作者贡献说明：}单人完成。研究领域：自然语言处理。

\vspace{0.5em}
\noindent\textbf{AI 使用声明：}本报告撰写过程中使用了 DeepSeek 与 Cursor AI 辅助组织语言与润色、表格设计与格式规范。
\newpage
\setcounter{page}{1}

\section{问题介绍}

随着智能家居与物联网技术的发展，家庭电力消耗的监测与预测在节能减排、降低用电成本以及智能电网调度中具有重要价值。家庭用电受季节变化、节假日、成员行为模式、设备类型及气象条件等多因素影响，呈现明显的非线性与周期性，准确预测具有较大挑战性和现实意义。

本实验基于 UCI 公开的 Individual Household Electric Power Consumption 数据集\cite{uci}。该数据来自法国 Sceaux 一户家庭，采样粒度为分钟，时间跨度为 2006 年 12 月至 2010 年 11 月。实验根据最近 90 天的多变量用电与气象信息，预测未来日级总有功功率，对应字段为 \varname{global_active_power}。短期任务预测未来 90 天，长期任务预测未来 365 天；两种任务分别训练独立模型，参数不共用。输入为过去 90 天的多变量序列。评价指标为均方误差 MSE 与平均绝对误差 MAE；每种配置重复 5 次实验，报告均值与标准差。

本实验采用 Direct Multi-step Forecasting 策略，即直接多步预测：模型在一次前向传播中输出完整的 90 或 365 日预测序列，而非逐步自回归滚动。该做法与短期、长期分别训练的要求一致，推理效率较高。

电力数据来自 UCI 仓库\cite{uci}，气象数据来自法国 Météo-France 月度气候数据集\cite{meteo}。课程未提供官方训练集与测试集划分文件，本实验按时间顺序划分：训练集为 2006-12-16 至 2008-12-31，共 747 天；测试集为 2009-01-01 至 2010-11-26，共 695 天。测试集长度满足长期预测所需的最少 455 天窗口。

原始数据为分钟级记录，缺失值以“?”标记，约占 1.25\%。按作业要求，全局有功功率、全局无功功率与三路分表能耗均按天求和；电压与全局电流强度按天求平均；气象变量 \varname{RR}、\varname{NBJRR1}、\varname{NBJRR5}、\varname{NBJRR10}、\varname{NBJBROU} 为月度数据，广播至该月各天。其中第三路分表对应气候控制系统。分钟级衍生特征 \varname{sub_metering_remainder} 按公式
\begin{equation}
\text{sub\_metering\_remainder} = \frac{\text{global\_active\_power} \times 1000}{60} - (\text{sub\_1} + \text{sub\_2} + \text{sub\_3})
\end{equation}
计算后按日求和。预处理阶段由日期字段生成星期、月份、年内序号与是否周末等日历特征，并与其他变量一并纳入模型输入。

气象站选用 Hauts-de-Seine 省 FONTENAY--CEN 站，编号 92032001。该站距 Sceaux 最近，且在 2006--2010 年有完整记录。\varname{RR} 原始值为毫米的十分之一，已除以 10 换算；\varname{NBJBROU} 在该站缺失，统一填 0。标准化在训练集上拟合，测试集仅做变换，避免信息泄漏。

样本采用步长为 1 的滑动窗口构造\cite{window1,window2,window3}。每个样本的输入为 $\mathbf{X} \in \mathbb{R}^{90 \times F}$，输出为 $\mathbf{y} \in \mathbb{R}^{H}$，其中 $H$ 取 90 或 365，特征维数 $F=17$。训练集含短期样本 568 个、长期 293 个；测试集含短期 516 个、长期 241 个。

\section{模型}

本实验实现 LSTM、Transformer 与 CNN-Transformer 三种模型。短期任务为 90 天输入、90 天输出；长期任务为 90 天输入、365 天输出。两类任务分别独立训练，基于 PyTorch 框架实现。

\subsection{LSTM 模型}

LSTM 模型以两层、隐层维度 128 的 LSTM 编码过去 90 天序列，取最后时刻隐状态，经两层全连接网络映射到 $H$ 步输出。给定输入 $\mathbf{x} \in \mathbb{R}^{90 \times F}$，有 $\mathbf{h}_T = \mathrm{LSTM}(\mathbf{x})$，$\hat{\mathbf{y}} = \mathrm{MLP}(\mathbf{h}_T)$，损失为 $\mathcal{L} = \mathrm{MSE}(\hat{\mathbf{y}}, \mathbf{y})$。LSTM 擅长建模序列依赖\cite{lstm}，作为经典基线。

\subsection{Transformer 模型}

Transformer 模型将输入线性投影至 $d_{\mathrm{model}}=64$，叠加正弦位置编码，经两层、四头注意力的 Encoder，在时间维均值池化后，再经 MLP 输出 $H$ 步。自注意力可建模长程依赖\cite{transformer}。Zeng 等\cite{zeng2023}指出标准 Transformer 在时序任务上对局部连续模式建模不足，这也是设计改进模型的出发点。

\subsection{CNN-Transformer 改进模型}

改进模型在 Transformer 前串联三层 Conv1D，卷积核依次为 3、5、7，再经位置编码、Encoder 与 MLP 一次性输出 $H$ 步。卷积与注意力均为常见组件，但具体组合与超参数针对本任务自行确定，并未复现某一已有网络。

设计过程如下。基线 Transformer 在本数据上短期表现尚可，但家庭用电存在明显日周期与周周期，纯注意力对短周期归纳偏置较弱\cite{zeng2023}，因此在编码前增加局部特征提取。输入为 90 天日粒度序列，自行选取 kernel 为 3、5、7 的三层卷积，分别约对应 3 天、1 周与 2 周感受野，由细到粗抽象局部模式。Kleeberg 等\cite{kleeberg2021}在同一 UCI 数据集上采用 CNN-LSTM 并融合气象数据；本模型改用 CNN-Transformer，并采用 Direct 多步输出，卷积核与层数均不同。相对 LSTM，本模型以注意力替代逐步递归并加入卷积前端；相对 Transformer，在编码前显式提取局部周期特征。

\subsection{训练设置}

优化器为 Adam，学习率 $10^{-3}$，损失为 MSE。短期 batch size 为 32，长期为 16；最大 epoch 80，早停 patience 12；随机种子取 42、123、456、789、2026。在训练集滑动窗口样本中按时间顺序取末尾 10\% 作验证集，用于早停。最终 MSE 与 MAE 在 2009 年后的独立测试集上计算，测试时段与训练时段不重叠。

\section{结果与分析}

\subsection{定量结果}

表~\ref{tab:results} 汇总三种模型在测试集上的结果。指标在反归一化后的原始尺度上计算，各配置重复 5 次实验，表中为均值与标准差。

\begin{table}[htbp]
    \centering
    \caption{各模型测试集 MSE 与 MAE}
    \label{tab:results}
    \begin{tabular}{llcc}
        \toprule
        模型 & 预测长度 & MSE & MAE \\
        \midrule
        LSTM            & 短期 90 天  & $297527.66 \pm 8633.30$  & $415.48 \pm 8.05$ \\
        LSTM            & 长期 365 天 & $\mathbf{222779.94 \pm 11713.97}$ & $\mathbf{353.23 \pm 10.95}$ \\
        Transformer     & 短期 90 天  & $262198.22 \pm 25517.02$ & $389.47 \pm 20.57$ \\
        Transformer     & 长期 365 天 & $290486.83 \pm 26657.28$ & $412.45 \pm 25.34$ \\
        CNN-Transformer & 短期 90 天  & $263430.64 \pm 69678.44$ & $385.40 \pm 56.65$ \\
        CNN-Transformer & 长期 365 天 & $226228.67 \pm 24559.67$ & $355.72 \pm 23.84$ \\
        \bottomrule
    \end{tabular}
\end{table}

\subsection{结果讨论}

长期预测任务上，LSTM 的 MSE 与 MAE 最低，且五次实验的标准差较小，MAE 标准差约为 10.95，训练较稳定。这可能因为 LSTM 的递归结构更适合平滑、缓慢的年度趋势，而一次性输出 365 步对 Transformer 类模型的容量与优化要求更高。短期预测上，Transformer 的 MSE 最低，为 262198，三种模型差距不大，说明在日尺度与周尺度规律下各架构均可学习。CNN-Transformer 短期 MSE 标准差为 69678，高于 LSTM 的 8633 与 Transformer 的 25517，对初始化较敏感；长期任务方差有所降低，但仍高于 LSTM。总体而言，LSTM 最稳定，Transformer 居中，CNN-Transformer 短期波动最大。

\subsection{预测曲线}

图~\ref{fig:short} 与图~\ref{fig:long} 分别给出短期与长期预测的三模型对比，图~\ref{fig:single} 为单模型示例。

\clearpage
\begin{figure}[H]
    \centering
    \fitfigure{../experiments/results/comparison_short.png}
    \caption{短期 90 天三种模型预测曲线与真实值对比}
    \label{fig:short}
\end{figure}

\clearpage
\begin{figure}[H]
    \centering
    \fitfigure{../experiments/results/comparison_long.png}
    \caption{长期 365 天三种模型预测曲线与真实值对比}
    \label{fig:long}
\end{figure}

\clearpage
\begin{figure}[H]
    \centering
    \begin{subfigure}[t]{0.48\textwidth}
        \centering
        \fitfigure{../experiments/results/lstm_short/seed_42/plot_seed42.png}
        \caption{LSTM 短期预测，seed 为 42}
    \end{subfigure}
    \hfill
    \begin{subfigure}[t]{0.48\textwidth}
        \centering
        \fitfigure{../experiments/results/cnn_transformer_long/seed_42/plot_seed42.png}
        \caption{CNN-Transformer 长期预测，seed 为 42}
    \end{subfigure}
    \caption{单模型预测示例}
    \label{fig:single}
\end{figure}

从曲线可见，模型能捕捉总体趋势与部分周期波动，但在用电峰值与突变日上仍有偏差。长期预测后半段误差有所增大，这是 Direct 策略的常见现象。长期任务中 LSTM 曲线整体更贴近真实值；CNN-Transformer 在部分区段对局部波动响应更明显。

\section{讨论}

本实验采用 Direct 策略，一次性输出 90 或 365 步预测，推理速度快，且与短期、长期分别训练的要求一致。其不足在于输出维数较高，365 步预测时对模型容量要求更大。此外，当前输入仅含过去 90 天的日历特征，预测未来各天时模型并不知道该日的星期与月份；而这些信息在实际部署中可以预先计算，后续可为每个预测步构造未来日历特征加以改进。

第三路分表能耗已按日求和。气象站因 Sceaux 本站停测而选用邻近 Fontenay 站；\varname{NBJBROU} 缺失处填 0。验证集为训练窗口末尾 10\%，仅用于早停；测试集为 2009 年后独立时段，与训练无重叠。

CNN-Transformer 在长期任务上接近 LSTM，MSE 分别为 226229 与 222780；短期 MAE 略优于 LSTM，分别为 385.40 与 415.48，但未全面超越基线，且不同 seed 间方差较大。文献\cite{zeng2023,kleeberg2021}用于说明问题背景与数据集参照，模型结构为自行设计。后续可引入未来日历特征，尝试 Encoder--Decoder 结构，并对 CNN-Transformer 增加学习率 warmup 以提升稳定性。

完整可运行代码已托管于 GitHub，链接见文首。克隆仓库后安装依赖，将 UCI 原始数据放入数据目录，依次运行天气下载、预处理与实验脚本即可复现。

\begin{thebibliography}{99}
\bibitem{uci}
UCI Machine Learning Repository.
\textit{Individual Household Electric Power Consumption}.
\url{https://archive.ics.uci.edu/dataset/235/individual+household+electric+power+consumption}

\bibitem{meteo}
data.gouv.fr.
\textit{Données climatologiques de base - mensuelles}.
\url{https://www.data.gouv.fr/fr/datasets/donnees-climatologiques-de-base-mensuelles}

\bibitem{lstm}
Hochreiter, S., \& Schmidhuber, J. (1997).
Long Short-Term Memory.
\textit{Neural Computation}, 9(8), 1735--1780.

\bibitem{transformer}
Vaswani, A., Shazeer, N., Parmar, N., et al. (2017).
Attention Is All You Need.
\textit{Advances in Neural Information Processing Systems}, 30.

\bibitem{zeng2023}
Zeng, A., Chen, M., Zhang, L., et al. (2023).
Are Transformers Effective for Time Series Forecasting?
\textit{Proceedings of the AAAI Conference on Artificial Intelligence}, 37(9), 11121--11128.

\bibitem{kleeberg2021}
Kleeberg, J., et al. (2021).
Energy Consumption Patterns and Load Forecasting with Profiled CNN-LSTM Networks.
\textit{Processes}, 9(11), 1870.

\bibitem{window1}
滑动窗口样本构造.
\url{https://blog.csdn.net/qq_47885795/article/details/143462299}

\bibitem{window2}
时间序列窗口划分.
\url{https://blog.csdn.net/weixin_39653948/article/details/105431099}

\bibitem{window3}
时间序列预测样本构造.
\url{https://datac.blog.csdn.net/article/details/105928752}

\end{thebibliography}

\end{document}
